\bigskip En 1992, le film \textit{Tous les matins du monde} reçoit le César de la meilleure musique originale. Ce prix consacre la performance du gambiste Jordi Savall et la découverte de la viole de gambe par le grand public. L’engouement suscité par cet instrument ancien illustre le renouveau de la musique ancienne en France, amorcé des décennies auparavant par les précurseurs de ce qui s'appellera par la suite la pratique historiquement informée de la musique.

\bigskip La pratique historiquement informée, aussi connue sous l’acronyme anglais HIP (\textit{Historically Informed Practice}), est un mouvement musical qui associe la recherche historique à l’interprétation d’oeuvres de «~musique ancienne~» pour recréer l’oeuvre au plus près de ses conditions de création.

\smallskip\indent Cette quête des sonorités du passé repose sur l’utilisation d’instruments d’époque (copies ou instruments restaurés), sur la redécouverte de partitions oubliées, sur l’analyse de traités d’harmonie, de jeu ou de facture instrumentale et sur des représentations iconographiques. Elle vise à redonner aux musiciens le geste et la matérialité du son du passé\footnote{Gétreau (Florence), Framboisier (Alban), His (Isabelle), \textit{Le son des musiques anciennes (1880-1980), Imaginer, fabriquer}, Rennes, Presses Universitaires de Rennes, 2024.} tout en nourrissant l’histoire des sensibilités et des paysages acoustiques\footnote{Chassagnette (Axelle), «~Entendre la musique du passé : ce que la pratique des instruments anciens peut apprendre aux historiens~», \textit{Les Carnets du LARHRA} [en ligne], 1 | 2017, mis en ligne en 2019 [réf. du 19 juin 2026], URL : \url{https://publications-prairial.fr/larhra/index.php?id=321}}.

\smallskip\indent Les chercheurs adeptes de la pratique historiquement informée s’attellent à théoriser la musique d’une époque et le style d’un compositeur à partir de traités et d’archives. Leurs découvertes et connaissances offrent aux musiciens, chefs et instrumentistes la matérialité perdue, pourtant nécessaire à la construction de leur interprétation musicale. Tous visent le même objectif : offrir au public une interprétation fidèle de l'œuvre.

\bigskip Est-il seulement possible de jouer la musique du passé comme elle était interprétée à l’époque ?

\smallskip\indent Si «~retrouver les musiques du passé est un point de mire très stimulant [...], personne ne peut dire honnêtement qu’on arrivera un jour à les faire sonner comme elles étaient\footnote{Phrase prononcée par Rémy Campos, professeur d'Histoire de la musique au CNSMDP dans l'émission Le Cours de l'Histoire. Référence : Mauduit, (Xavier), animateur. Musiques anciennes, et je recrée le son ! Avec Rémy Campos et Florence Gétreau. [Série de podcasts audio] Histoire et musique, l'accord parfait, France Culture, 2025. \url{https://www.radiofrance.fr/franceculture/podcasts/le-cours-de-l-histoire/musiques-anciennes-et-je-recree-le-son-5967778}}~». 

\smallskip\indent Cette recherche d'authenticité se heurte d'abord à la polysémie de la notion d'«~œuvre originale~». Que recouvre vraiment cette notion ? Prenons l'exemple de la célèbre œuvre de Rameau, \textit{Les Indes galantes}. L'œuvre originale correspond-elle au manuscrit autographe perdu, aux premières copies contemporaines de l'œuvre indiquant la manière de la jouer, ou bien à l’interprétation donnée lors de la première représentation au théâtre du Palais-Royal, le 23 août 1735 ?

\smallskip\indent Aucune de ces pistes ne suffit à définir l'original authentique et immuable de l'œuvre musicale. En effet, réduire l'œuvre à son manuscrit autographe poserait le problème de la notation ancienne, souvent partielle ou codifiée, ainsi que des enjeux de sa compréhension et de sa transmission. S'en tenir aux premières copies avec les indications d'époque est tout aussi problématique car cela reviendrait à occulter l'importance de l'improvisation et de l'ornementation, pourtant fondamentales dans la musique baroque. Cependant, assimiler l'œuvre au son de la première création relève également de l'illusion en raison de l'absence de tout outil d'enregistrement sonore avant la mise au point du phonographe, à la fin du XIX$^{e}$ siècle.

\bigskip Le besoin d’authenticité de l'œuvre musicale, par la documentation de sa genèse, tient donc autant d’un intérêt patrimonial que musical. Il convient de rapprocher la quête des pionniers de la \enquote{résurrection de la musique du passé}\footnote{Landowska (Wanda), Musique ancienne, Paris, 1909} du XIX$^{e}$ siècle, du contexte de la naissance de la notion contemporaine de patrimoine. L’ère de la vague romantique, en vogue au XIX$^{e}$ siècle, encourage artistes et politiques à porter davantage attention au patrimoine historique, notion encore en construction à l’époque.

\smallskip\indent Ce dernier négligé et parfois même vandalisé, au grand dam de Victor Hugo qui s’en offusque dans \textit{Guerre aux démolisseurs}\footnote{\enquote{\textit{Quoique appauvrie par les dévastations révolutionnaires, par les spéculateurs mercantiles et surtout par les restaurateurs classiques, la France est riche encore en monuments français. Il faut arrêter le marteau qui mutile la face du pays. Une loi suffirait, qu’on la fasse […] Il y a deux choses dans un monument : son usage et sa beauté. Son usage appartient au propriétaire, sa beauté à tout le monde. C’est donc dépasser son droit que de le détruire}} (Victor Hugo, \textit{Guerre aux démolisseurs}, 1825).}, fait l’objet, dans une visée mémorielle, d’une volonté de sauvegarde dépassant les clivages politiques hérités de la Révolution. La création en 1830 de l’Inspection générale des Monuments historiques constitue un jalon institutionnel important de la sauvegarde du patrimoine.
En musique, comme pour le patrimoine monumental, le XIX$^{e}$ siècle consacre la volonté de sauvegarder les oeuvres du passé\footnote{\enquote{\textit{Au XIX$^{e}$ siècle, le public a cessé de croire que la musique de son temps était un absolu ; il a découvert une valeur artistique dans des œuvres anciennes dans la mesure même où la musique de ses contemporains cessait souvent de lui donner satisfaction : dès les derniers quatuors de Beethoven, Berlioz, Wagner, le langage des musiques nouvelles a souvent semblé d'un accès difficile et le commun des auditeurs s'est réfugié dans la sécurité des musiques du passé}}. Wangermée (Robert), La redécouverte des musiques du passé et la quête de l'authenticité. Expériences bruxelloises. In: \textit{Bulletin de la Classe des Beaux-Arts}, tome 9, n°7-12, 1998. pp. 157-179. DOI : \url{https://doi.org/10.3406/barb.1998.20506}}.

\smallskip\indent La conscience patrimoniale serait donc, selon la conception de ses pionniers, l’intérêt majeur de la \enquote{musique ancienne}. À cette époque, la musique était avant tout une musique du présent qui occultait les répertoires plus anciens, ceux antérieurs aux périodes classique et romantique, c’est-à-dire les périodes médiévale, Renaissance et baroque. Des pionniers, comme Geneviève de Chambure, Wanda Landowska ou encore Henri Casadesus, prônaient alors un retour aux instruments anciens et la redécouverte de sources musicales anciennes\footnote{L'histoire des pionniers de la \enquote{musique ancienne} sera abordée plus en détails dans le chapitre 2 }.

\bigskip Devant cette quête difficile, voire impossible du son original de l'œuvre, se pose la question de l’intérêt de la pratique historiquement informée.

\smallskip\indent L’intérêt de la pratique de la musique ancienne, s’il n’est pas dans l’authenticité de l'œuvre, réside donc dans le processus d’interprétation de l'œuvre. Avec le bouleversement de la musique ancienne dans les années 1970, la naissance de la pratique historiquement informée opère également un changement dans le processus d’interprétation.

\smallskip\indent L’intérêt de la musique ancienne --- qui résidait alors dans la découverte, l'édition et la diffusion des partitions anciennes --- devient dès lors un projet artistique et scientifique. Ce projet couronne un processus intellectuel de travail historique sur les sources, enrichissant l’édition de partitions de recherches historiques sur le contexte de l'œuvre. Il passe également par une attention portée sur les instruments anciens. Les chercheurs et les luthiers s’unissent donc pour répliquer ou restaurer des instruments. Ainsi, la pratique historiquement informée se situe entre l’archéologie expérimentale, les \textit{performance studies} et la recherche historique et musicologique. Elle mobilise des disciplines phares telles que la facture instrumentale, l’organologie, la paléographie ou encore l’ecdotique musicale. La perspective des \textit{performance studies} prend alors tout son sens : l'œuvre n'est plus un texte figé (le manuscrit original, source musicale), mais un événement joué, vivant et éphémère (sur la scène).

\smallskip\indent C'est tout le paradoxe de la démarche. Alors que le public cherchait dans le passé la sécurité de répertoires connus face aux musiques nouvelles comme le constatait Robert Wangermée\footnote{Wangermée (Robert), La redécouverte des musiques du passé et la quête de l'authenticité. Expériences bruxelloises. In: \textit{Bulletin de la Classe des Beaux-Arts}, tome 9, n°7-12, 1998. pp. 157-179. DOI : \url{https://doi.org/10.3406/barb.1998.20506}}, la pratique historiquement informée provoque de l'inattendu et de la surprise au cœur de ce patrimoine.

\smallskip\indent L’horizon d’attente des mélomanes change alors, au moment où la musique cesse d’être un art du présent. Le public préoccupé par la question du patrimoine musical s’intéresse aux musiques anciennes et pense retrouver, dans la musique d’autrefois, une pureté musicale.

\bigskip Toutefois, l’idéalisation de la musique du passé se heurte là encore à la réalité de la musique historique. En effet, les contraintes techniques liées aux instruments anciens --- à l’image des cordes en boyau qui, plus sensibles aux conditions d’hygrométrie, se désaccordent plus vite --- en complexifient l’interprétation et l’écoute. L’exemple du diapason est une illustration marquante des particularités de la musique ancienne. Le diapason, variable selon les époques et les localités, est un réel enjeu de transmission musicale. Dans la musique ancienne, dite baroque, des XVII$^{e}$ et XVIII$^{e}$ siècles, le diapason était fixé à 415 hertz (baroque allemand), alors qu’il est aujourd’hui de 440 hertz tel que fixé à la conférence internationale de Londres en 1953. On ne s'étonnera alors pas que l’horizon d’attente d’un auditeur du XXI$^{e}$ siècle soit quelque peu dérouté par une pièce baroque.

\bigskip Les années 1950-1970 sont le théâtre d’une accélération de la \enquote{musique ancienne}. Comme nous l’avons vu, ce qui n’était alors qu’une volonté de connaître le répertoire, motivée peut-être par un manque de reconnaissance de la musique contemporaine, devient un projet artistique à part entière, ouvrant la voie à la création d’ensembles majeurs de la scène de ce qui s’appelait alors la musique baroque. Le Concentus Musicus Wien créé par Nikolaus Harnoncourt, Les Arts Florissants créés par William Christie, ou encore la Chapelle Royale créée par Philippe Herreweghe sont les plus connus d'entre eux. La notoriété du film \textit{Tous les matins du monde} (1992) a renforcé cet engouement et aujourd’hui encore, vingt-quatre ans après la popularisation du terme de pratique historiquement informée par l’ouvrage de John Butt \textit{Playing with History}, ce mouvement musical occupe encore le devant de la scène.

\bigskip Pour répondre à l’exigence de la musique historiquement informée, les musiciens, chefs et chercheurs sont donc à la recherche des sources de la musique ancienne. Partitions manuscrites autographes, traités d'instruments et d'interprétation ou tout autre document permettant de documenter les œuvres sont autant de sources recherchées par les musiciens en ce qu'elles leur permettent de retrouver le geste et l'esprit de la musique inscrite sur leur partition. Loin d'être un ensemble de documents de musique notée prenant la poussière sur une étagère, le patrimoine des bibliothèques musicales et d’orchestre est au contraire témoin des attentes du compositeur et moteur de la créativité de l'interprétation historiquement informée.

\smallskip\indent Dès lors, les bibliothèques d'orchestre, parfois vues comme de simples conservatoires de partitions, deviennent au contraire des centres de ressources dont la visée première est de nourrir l'interprétation des musiciens au moyen d’un lien fort avec la recherche. Outre les missions traditionnelles des bibliothèques d'orchestre --- sélection des éditions des oeuvres à jouer, location ou achat du matériel, préparation du matériel (coups d'archet, indications du chef...) --- celles-ci se voient confier un nouveau rôle qui est de documenter les œuvres et de patrimonialiser les partis pris des chefs et des musiciens.

\bigskip Le patrimoine des bibliothèques musicales est donc un patrimoine singulier constitué à la fois de documents patrimoniaux (comme le manuscrit autographe du \textit{Don Giovanni} de Mozart conservé à la Bibliothèque nationale de France) et d'\enquote{archives intermédiaires} (comme des conducteurs de symphonies de Beethoven annotées de la main de la cheffe Laurence Equilbey, à la bibliothèque de l’orchestre Insula Orchestra, qui sont encore utilisés pour les productions).

Dans une bibliothèque d’orchestre, cette dualité de typologie de documents, auxquels il faut ajouter bien sûr les usuels sur la musique et les archives courantes de l’orchestre, pose la question du statut des collections\footnote{Hausfater (Dominique) et Campos (Rémy), L’interprétation musicale dans les fonds des bibliothèques : Journée d’étude Du 7 Janvier 2006 Conservatoire de Musique de Genève/Haute École de Musique. \textit{Fontes Artis Musicae}, vol. 54, no. 1, 2007, pages 1–5. JSTOR, \url{http://www.jstor.org/stable/23510565} (consulté le 25 août 2026).}.
Quel peut être le statut des collections d’une bibliothèque d’orchestre alors même que la nature de ce type de bibliothèque défie la notion du cycle de vie des documents, en patrimonialisant des archives d'exécution encore en usage au cœur même de l'activité scénique ?\footnote{Cette situation hybride interroge la théorie des trois âges des archives (formalisée par Yves Pérotin). En défiant la séparation étanche entre âge intermédiaire et âge définitif, le matériel d'orchestre annoté conserve simultanément sa valeur primaire d'usage fonctionnel pour la pratique scénique et sa valeur secondaire de preuve mémorielle, des choix des chefs par exemple. Il constitue ainsi une « archive d'exécution », témoin matériel fondamental pour la recherche musicologique et l'histoire de l'interprétation.}

\bigskip Les promesses du numérique offrent en effet aujourd'hui une nouvelle ère à la pratique historiquement informée. À l'heure du numérique, la recherche en musique historique bénéficie de nouveaux outils pour aboutir sa quête du geste et du son. Dans un contexte de transition numérique, où \enquote{les bibliothèques [se trouvent] face au monde des données}\footnote{Mesguich (Véronique), \textit{Les bibliothèques face au monde des données}, Presses de l’enssib, 2023, \url{https://doi.org/10.4000/books.pressesenssib.17686}.}, l'étude du potentiel offert par les technologies numériques, liant bibliothèques et musiciens dans une réflexion sur la médiation des sources patrimoniales, est primordial. 

\smallskip\indent Si les pionniers de la pratique historiquement informée, qui n’était alors que la «~musique ancienne~», souhaitaient simplement jouer la musique du passé en redécouvrant des pièces oubliées, les chercheurs et musiciens d’aujourd’hui rêvent de jouer la musique ancienne comme elle était jouée à l’époque.

\bigskip C’est précisément à ces questionnements que tente de répondre la nouvelle plateforme LaDocumenta.eu, mise au point par un consortium d’orchestres européens autour de l’orchestre \textit{Insula Orchestra}. Cette plateforme souhaite en effet conserver le patrimoine musical de la pratique historiquement informée tout en n’occultant pas son caractère singulier : un patrimoine nécessitant d’être conservé pour les siècles futurs mais également toujours en utilisation aujourd’hui pour des projets musicaux.

\bigskip La construction d’un tel projet illustre à quel point le numérique peut être une ressource pour la pratique historiquement informée. La notion d'interopérabilité des contenus musicaux, dont les documents comportent à la fois de la musique notée et du texte, est l'exemple type des difficultés éprouvées quant à la visibilité des documents.

\smallskip\indent Si la numérisation des partitions annotées, essentielle pour conserver la mémoire des choix artistiques des chefs, permet de rendre accessible nombre de ressources musicales, le chemin est encore long pour permettre l'exploitation de ces ressources. L'enjeu d'une telle plateforme est d'opérer un changement d'échelle de la ressource musicale pour bénéficier des évolutions bibliothéconomiques dont la réflexion est engagée dans le projet de la Transition Bibliographique.

\smallskip\indent En effet, afin de rendre les ressources musicales historiquement informées interopérables, les documents physiques conservés dans les bibliothèques du consortium doivent devenir, dans le cadre de la plateforme, des données numériques structurées et atomiques capables de s'inscrire dans le web sémantique et de garantir ainsi leur visibilité\footnote{La structuration et l'interopérabilité des sources musicales et de leurs usages scéniques s'inscrivent dans l'écosystème des standards internationaux du web de données et du patrimoine écrit (notamment les modèles conceptuels du modèle IFLA-LRM / LRMoo pour le spectacle vivant, et les schémas d'encodage comme MEI pour la musique notée), dont la mise en œuvre et les limites seront analysées au cours de ce travail.}.

\bigskip LaDocumenta.eu n'est pas un projet de recherche sur la musique historiquement informée. C'est un projet qui étudie le rôle des technologies numériques dans la recherche musicologique et la quête de l'interprétation musicale historiquement informée. Elle est également le lieu d'inscription de telles institutions dans un réseau professionnel riche où la bibliothèque n'est pas refermée sur elle-même mais s'inscrit au coeur d'un réseau de bibliothèques au service de la recherche. Le partage des sources, dans une perspective de science ouverte pour les recherches musicologiques, est aussi à destination des \textit{performance studies} de la musique historiquement informée.

\bigskip Ce mémoire s'inscrit à la croisée des thématiques de la musicologie, de la bibliothéconomie et des humanités numériques. \enquote{De la source à la scène} est un titre inspiré de mon expérience de stage de fin d'études en tant que bibliothécaire-musical assistant à l'orchestre \textit{Insula Orchestra}. Il vise à souligner l’aspect continuel de la pratique historiquement informée, du travail préparatoire en bibliothèque, sur les sources patrimoniales musicales, jusqu’à la scène qui est le lieu de l’aboutissement des recherches sur l’interprétation.

\bigskip Le coeur de cette expérience à la Seine Musicale (Boulogne), lieu de résidence de cet orchestre de pratique historiquement informée, a été de prendre part au projet LaDocumenta.eu. Prenant la forme d'une bibliothèque numérique, cette plateforme se veut être une aide à destination des orchestres pour organiser la gestion de leur bibliothèque d'orchestre. L'enjeu est également de garantir la patrimonialisation de la musique historiquement informée et d'oeuvrer pour sa reconnaissance comme patrimoine immatériel de l'UNESCO. Aujourd'hui soutenue par des fonds européens, cette plateforme est à la fois un outil de médiation de contenus culturels comme Europeana et un outil de gestion de bibliothèques pouvant remplacer des logiciels métiers comme PMB\footnote{PMB est un logiciel libre de gestion de bibliothèques (SIGB), il est l'outil de référence pour les bibliothèques musicales et d'orchestre, utilisé notamment chez Radio France, au Centre Musique Baroque de Versailles (CMBV), ou encore dans les bibliothèques des Conservatoires Nationaux Supérieurs de Paris et Lyon (CNSMD).}. Elle souligne les problématiques institutionnelles de la bibliothéconomie musicale en contexte numérique.

\bigskip Or ces problématiques institutionnelles se doublent de problématiques concrètes - le dépôt des ressources, l'enrichissement des notices et la médiation des collections - que les moyens traditionnels peinent à satisfaire et qui sont des terrains d'expérimentation pour les nouvelles technologies. L'IA est alors explorée comme une hypothèse d'assistance à la description archivistique et au traitement de masse, dont il convient d'évaluer la pertinence scientifique et la reproductibilité.

\smallskip\indent Les programmes de financements européens le montrent en proposant des sujets très orientés vers ces questionnements : l'appel à projet \enquote{Soutien à la découvrabilité des contenus culturels francophones dans l’environnement numérique à l’ère de l’IA} de la stratégie France-Québec 2025-2030, l'appel à projet ANR PEPR-ICCARE \enquote{Industries culturelles et créatives}, ou encore les appels Europe creative (dont LaDocumenta.eu a été lauréat) et Horizon Europe sont autant d'exemples qui montrent que la recherche dans le domaine culturel est attendue sur ces thématiques.

\bigskip À l'aune de ces réflexions, il convient de se demander : en quoi les plateformes numériques musicales offrent-elles une solution originale pour transformer le patrimoine des bibliothèques d'orchestre en données de recherche exploitables pour construire une interprétation de pratique historiquement informée ?

\bigskip L'étude de ce projet, de la source patrimoniale à la scène, nous amène à ancrer dans un premier temps la bibliothèque d'orchestre dans son acception institutionnelle européenne, entre la conservation patrimoniale, l'interprétation scénique et les nouvelles technologies numériques. Dans un second temps, cette mise en contexte explique la transformation technique et intellectuelle du document. Transformer la partition annotée ou encore les traités d'interprétation en données structurées suppose en effet de repenser la politique documentaire sous le prisme de l'interopérabilité et des enjeux du web sémantique. Dans un troisième et dernier temps, nous verrons que le recours à l'IA ne peut s'envisager qu'après cette transformation, l'IA nécessitant des données structurées et exploitables. Dès lors l'IA peut remplir un rôle d'assistant documentaire dont il convient de circonscrire les usages et les limites tant éthiques que méthodologiques au service de la recherche et de l'interprétation historiquement informée.